\songtitle{Límite sincopado}

\subsection{Notes}

\subsection{Style}


\begin{lyrics}
\songpart{Intro}
\annotation{Spoken, piano in the background}

Él se frustra porque no logra lo que quiere.\\
Tú te frustras porque no hace lo que se debe.

\annotation{Brass attack}\\
\annotation{Percussion enters}

\songpart{Verse 1}

No repitas la lección,\\
que ya conoce el deber;\\
no aumentes la presión,\\
si le cuesta proceder.\\
No juzgues por lentitud,\\
ni por tardar en responder;\\
a veces falta virtud\\
al medio para poder.

\songpart{Pre-Chorus}

No es desdén ni dejadez,\\
sino un límite tal vez.

\songpart{Chorus}

Es que uno se desespera,\\
y es frustrante no entender,\\
por qué falla la lección;\\
como si él no comprendiera\\
lo que tiene que hacer,\\
mas no es esa la razón.

\songpart{Verse 2}

No digas: "debes seguir",\\
ni "ponte a organizar";\\
procura más bien partir\\
la senda en pasos al andar.\\
Menos opciones darás,\\
más fácil será escoger;\\
si el peso logras quitar,\\
más libre podrá mover.

\songpart{Pre-Chorus}

Haz ligera la labor,\\
no le añadas más temor.

\songpart{Chorus}

Es que uno se desespera,\\
y es frustrante no entender,\\
por qué falla la lección;\\
como si él no comprendiera\\
lo que tiene que hacer,\\
mas no es esa la razón.

\songpart{Bridge}

\annotation{Spoken}

No es falta de capacidad.\\
No es falta de entendimiento.

El problema no es saber qué debe hacerse.\\
El problema es convertir ese entendimiento en acción.

A veces por sobrecarga.\\
A veces por agotamiento.\\
A veces porque el primer paso pesa más de lo que parece desde afuera.

Y cuando la ayuda nace de la frustración,\\
sin querer puede volver más pesado ese primer paso.

\songpart{Montuno}

\annotation{Call and Response}

\annotation{Group}\\
(Pero uno se desespera)

\annotation{Lead}\\
Es frustrante no saber

\annotation{Group}\\
(¿Por qué es que no le importa?)

\annotation{Lead}\\
No es desdén ni dejadez

\annotation{Group}\\
(Entonces ¿qué es lo que pasa?)

\annotation{Lead}\\
No le cuesta comprender

\annotation{Group}\\
(Entonces ¿qué es lo que pasa?)

\annotation{Lead}\\
Le cuesta poder hacer

\annotation{Group}\\
(¿Qué es lo que puedo hacer?)

\annotation{Lead}\\
Encamina y haz ligera la labor

\annotation{Mambo}

\annotation{Brass section}\\
\annotation{Piano montuno}\\
\annotation{Percussion break}

\songpart{Verse 3}

No tomes por desamor\\
lo que es dificultad real;\\
separa falta y dolor\\
de un límite funcional.\\
Pregunta qué obstáculo\\
le impide hoy avanzar;\\
y ofrece apoyo práctico,\\
sin reprochar ni empujar.

\songpart{Pre-Chorus}

Antes de juzgar su acción,\\
busca la causa y la razón.

\songpart{Final Chorus}

Es que uno se desespera,\\
y es frustrante no entender,\\
por qué falla la lección;\\
como si él no comprendiera\\
lo que tiene que hacer,\\
mas no es esa la razón.

\songpart{Outro}

No es desdén ni dejadez,\\
sino un límite tal vez.

Haz ligera la labor,\\
no le añadas más temor.

Antes de juzgar su acción,\\
busca la causa y la razón.

\annotation{Spoken}

A veces ayudar\\
es quitar peso del camino.

\annotation{Final brass hit}
\end{lyrics}

